\documentclass[12pt]{article}
\usepackage[utf8]{inputenc}
\usepackage{amsmath,amssymb,amsthm}
\usepackage[margin=1in]{geometry}

\newtheorem{theorem}{Theorem}
\newtheorem{lemma}[theorem]{Lemma}

\title{Problem 219: Skew-cost Coding}
\author{Project Euler Solution}
\date{}

\begin{document}
\maketitle

\section{Problem Statement}

Build a prefix-free binary code for $N = 10^9$ symbols. A \texttt{0} bit costs 1, a \texttt{1} bit costs 4. Minimize the total cost (sum of all codeword costs).

\section{Mathematical Foundation}

\begin{theorem}[Greedy optimality]
The minimum total cost is achieved by repeatedly splitting the minimum-cost leaf. Splitting a leaf of cost $c$ removes it and creates leaves of costs $c+1$ and $c+4$.
\end{theorem}

\begin{proof}
Suppose an optimal code has a non-split leaf $\ell_1$ of cost $c_1$ and a split leaf $\ell_2$ of cost $c_2 > c_1$. Swapping (making $\ell_2$ a leaf and splitting $\ell_1$ instead) changes the total cost by $c_1 - c_2 < 0$, contradicting optimality. By induction on the number of splits ($N-1$ total, starting from one leaf), the greedy strategy of always splitting the cheapest leaf is optimal.
\end{proof}

\begin{lemma}[Cost accounting]
Each split of a leaf with cost $c$ increases total cost by $c + 5$ and leaf count by $1$.
\end{lemma}

\begin{proof}
Removing cost $c$ and adding $c+1$ and $c+4$: net cost change is $(c{+}1)+(c{+}4)-c = c+5$. Net leaf change: $2-1=1$.
\end{proof}

\begin{theorem}[Total cost formula]
After $N-1$ splits starting from a single leaf of cost~0:
\[
\mathrm{Total} = 5(N-1) + \sum_{i=0}^{N-2} c_i,
\]
where $c_i$ is the cost of the leaf split at step $i$.
\end{theorem}

\begin{proof}
Initial cost is 0. By the lemma, split $i$ adds $c_i + 5$. Summing over all splits: $\mathrm{Total} = \sum_{i=0}^{N-2}(c_i+5) = 5(N{-}1) + \sum c_i$.
\end{proof}

\begin{theorem}[Bulk simulation]
If all $k$ leaves at the current minimum cost $c$ are split simultaneously, the result is identical to $k$ sequential greedy splits. Total cost increases by $k(c+5)$.
\end{theorem}

\begin{proof}
All $k$ leaves share the minimum cost $c$, so they are the next $k$ candidates for the greedy algorithm. Since the splits operate on distinct leaves and produce independent children, the order is immaterial, and the bulk update (remove $k$ at cost $c$; add $k$ at $c{+}1$ and $k$ at $c{+}4$) is correct.
\end{proof}

\section{Algorithm}

\begin{enumerate}
\item Maintain a histogram mapping each cost level to its leaf count.
\item Initialize: one leaf at cost 0, total $= 0$, leaf count $= 1$.
\item While leaf count $< N$:
  \begin{enumerate}
  \item Let $c$ be the minimum cost, $k$ its count (capped at $N - \text{leaves}$).
  \item Update: total $\mathrel{+}= k(c+5)$, leaves $\mathrel{+}= k$.
  \item Histogram: remove $k$ from cost $c$; add $k$ to costs $c{+}1$ and $c{+}4$.
  \end{enumerate}
\item Return total.
\end{enumerate}

\section{Complexity Analysis}

\begin{itemize}
\item \textbf{Time:} The minimum cost advances by at least 1 per iteration, and the maximum cost is $O(\sqrt{N})$ (leaf count grows super-linearly with cost). Total: $O(\sqrt{N}\log N)$.
\item \textbf{Space:} $O(\sqrt{N})$ distinct cost levels in the histogram.
\end{itemize}

\section{Answer}

\[
\boxed{64564225042}
\]

\end{document}
